\section{Non-Uniform Binary Fuse Filters}

This project was completed for CSC 592: Algorithms for Big Data.

\subsection{Motivation}

Distributed systems often require fast and small set data structures. For example, content delivery networks avoid caching one-hit wonders by querying a visited
set~\cite{maggs2015}. Supporting such exact membership queries on a large universe\footnote{Set of all keys.} $U$ requires $\Omega(\lg |U|)$ bits per key. For 
large enough sets, this is impractical. By permitting false positive results with some probability $\epsilon$, the lower bound is reduced to $\lg \frac{1}{\epsilon}$ 
bits per key~\cite{carter1978}. In other words, \textit{approximate} membership queries can be much more memory efficient than exact membership queries. \\

\noindent
A binary fuse filter~\cite{graf2022} is a recent static probabilistic data structure for approximate membership queries. Membership information is stored in an 
array of $m$-bit values partitioned into disjoint segments. To query the filter, a key is hashed to a fingerprint value and to three array entries located in 
consecutive segments (see Figure~\ref{fig:csc592_bigdata_bff}). If the bitwise-XOR of the entries is equal to the fingerprint, then the query \emph{might} be in 
the set. Otherwise, the query \emph{definitely} is not in the set.

\begin{figure}[H]
    \centering
    \begin{tikzpicture}[>=Stealth, arraynode/.style={draw, minimum width=7mm, minimum height=6mm, outer sep=0, inner sep=0}]

        \node[arraynode, fill=lightyellow] (a1) {000};
        \node[arraynode, fill=lightorange, right=0 of a1] (a2) {110};
        \node[arraynode, fill=lightorange, right=0 of a2] (a3) {000};
        \node[arraynode, fill=lightyellow, right=0 of a3] (a4) {000};
        \node[arraynode, fill=lightyellow, right=0 of a4] (a5) {000};
        \node[arraynode, right=0 of a5, path picture={
        \shade[left color=lightorange, right color=lightred]
                (path picture bounding box.south west)
                rectangle (path picture bounding box.north east);
        }] (a6) {010};
        \node[arraynode, right=0 of a6, path picture={
        \shade[left color=lightred, right color=lightpink]
                (path picture bounding box.south west)
                rectangle (path picture bounding box.north east);
        }] (a7) {001};
        \node[arraynode, fill=lightyellow, right=0 of a7] (a8) {000};
        \node[arraynode, fill=lightyellow, right=0 of a8, path picture={
        \shade[left color=lightred, right color=lightpink]
                (path picture bounding box.south west)
                rectangle (path picture bounding box.north east);
        }] (a9) {000};
        \node[arraynode, fill=lightyellow, right=0 of a9] (a10) {000};
        \node[arraynode, fill=lightyellow, right=0 of a10, path picture={
        \shade[left color=lightpink, right color=lightpurple]
                (path picture bounding box.south west)
                rectangle (path picture bounding box.north east);
        }] (a11) {111};
        \node[arraynode, fill=lightyellow, right=0 of a11] (a12) {000};
        \node[arraynode, fill=lightpurple, right=0 of a12] (a13) {000};
        \node[arraynode, fill=lightyellow, right=0 of a13] (a14) {000};
        \node[arraynode, fill=lightyellow, right=0 of a14] (a15) {000};
        \node[arraynode, fill=lightpurple, right=0 of a15] (a16) {000};

        \draw[thick] 
            ([yshift=4pt]a1.north west) -- 
            ([yshift=-4pt]a1.south west);
        
        \foreach \i in {2,4,6,8,10,12,14,16} {
            \draw[thick] 
                ([yshift=4pt]a\i.north east) -- 
                ([yshift=-4pt]a\i.south east);
        }

        \node[circle, draw, above=12mm of a4, fill=lightorange] (x) {$x$};
        \draw[->] (x) to (a2.north);
        \draw[->] (x) to (a3.north);
        \draw[->] (x) to (a6.north);
        
        \node[circle, draw, above=12mm of a7, fill=lightred] (y) {$y$};
        \draw[->] (y) to (a6.north);
        \draw[->] (y) to (a7.north);
        \draw[->] (y) to (a9.north);

        \node[circle, draw, above=12mm of a10, fill=lightpink] (w) {$w$};
        \draw[->] (w) to (a7.north);
        \draw[->] (w) to (a9.north);
        \draw[->] (w) to (a11.north);
        
        \node[circle, draw, above=12mm of a13, fill=lightpurple] (z) {$z$};
        \draw[->] (z) to (a11.north);
        \draw[->] (z) to (a13.north);
        \draw[->] (z) to (a16.north);
        
    \end{tikzpicture}
    \caption{Binary fuse filter with 16-entry array, 2-entry segments, and 3-bit fingerprints. Vertical bars mark segment boundaries. The filter contains keys 
    $\mathcolorbox{lightorange}{x}$, $\mathcolorbox{lightred}{y}$, $\mathcolorbox{lightpink}{w}$ and $\mathcolorbox{lightpurple}{z}$, with fingerprints 
    $\mathcolorbox{lightorange}{100}$, $\mathcolorbox{lightred}{011}$, $\mathcolorbox{lightpink}{110}$ and $\mathcolorbox{lightpurple}{111}$, respectively. 
    Entries hashed to by keys are marked with arrows and the keys' colors. Each key hashes to three consecutive segments.}
    \label{fig:csc592_bigdata_bff}
\end{figure}


\noindent
To construct a binary fuse filter, choose each array entry such that, for each key, the bitwise-XOR of the three corresponding array entries is equal to
the fingerprint. This requires solving a system of linear equations in $\mathbb{F}^{m}_2$. Let $h_1$, $h_2$ and $h_3$ be the hash functions, $f$ be the 
fingerprint function, and $A$ be the array. Each key $k$ yields a linear equation $A[h_1(k)] \oplus A[h_2(k)] \oplus A[h_3(k)] = f(k)$ that must be satisfied 
(see Figure~\ref{fig:csc592_bigdata_soe}). If some equation is not satisfied, the filter will incorrectly report the corresponding key is \emph{definitely} not 
in the set. Since false negatives are not permitted, the choice of array entries must satisfy every key's equation.

\begin{figure}[H]
\centering  
\[
\setlength{\arraycolsep}{2pt}
\begin{array}{cccccccccccccc}

\mathcolorbox{lightorange}{A[1]} 
&\oplus& \mathcolorbox{lightorange}{A[2]} 
&\oplus& \gradientmathbox{lightorange}{lightred}{A[5]} 
&=& \mathcolorbox{lightorange}{100} &\qquad

\gradientmathbox{lightorange}{lightred}{A[5]} 
&\oplus& \gradientmathbox{lightred}{lightpink}{A[6]} 
&\oplus& \gradientmathbox{lightred}{lightpink}{A[8]} 
&=& \mathcolorbox{lightred}{011}\\[1em]

\gradientmathbox{lightred}{lightpink}{A[6]} 
&\oplus& \gradientmathbox{lightred}{lightpink}{A[8]} 
&\oplus& \gradientmathbox{lightpink}{lightpurple}{A[10]} 
&=& \mathcolorbox{lightpink}{110} &\qquad

\gradientmathbox{lightpink}{lightpurple}{A[10]} 
&\oplus& \mathcolorbox{lightpurple}{A[12]} 
&\oplus& \mathcolorbox{lightpurple}{A[15]} 
&=& \mathcolorbox{lightpurple}{111}
\end{array}
\]
\caption{System of linear equations in $\mathbb{F}^3_2$ induced by Figure~\ref{fig:csc592_bigdata_bff}. Each $A[i]$ is the $i$-th entry in the backing array $A$
(zero-indexed). The keys are $\mathcolorbox{lightorange}{x}$, $\mathcolorbox{lightred}{y}$, $\mathcolorbox{lightpink}{w}$ and $\mathcolorbox{lightpurple}{z}$ 
with fingerprints $\mathcolorbox{lightorange}{100}$, $\mathcolorbox{lightred}{011}$, $\mathcolorbox{lightpink}{110}$ and $\mathcolorbox{lightpurple}{111}$, 
respectively. Array entries and fingerprints are marked with corresponding key color(s). There is one equation for each key.}
\label{fig:csc592_bigdata_soe}
\end{figure}

\noindent
To solve the system, model it as a hypergraph. A hypergraph is a graph where edges may contain any subset of vertices. For each entry, create a vertex. For each 
equation, create an edge containing entries (vertices) in the equation (see Figure~\ref{fig:csc592_bigdata_hg}). Then, 2-peel the hypergraph, that is:
\begin{enumerate}
    \item Remove vertices of degree zero, then set their entries to $0^m$.
    \item Remove vertices of degree one and their incident edges. Then, set their entries to the bitwise-XOR of the fingerprint and the two other entries in the
    edge. If the other entries are not yet set, defer until they are set.
    \item Repeat until there are no vertices of degree less than two, or until the empty hypergraph is reached (i.e. all entries set).
\end{enumerate}
If the empty hypergraph is reached, the hypergraph is called 2-peelable. The system is solved if and only if its hypergraph is 2-peelable. \\

\begin{figure}[H]
\centering
\begin{tikzpicture}[scale=1.2]

\node (v1) at (0, 2) {};
\node (v2) at (0, 1) {};
\node (v3) at (2, 1.5) {};
\node (v4) at (4, 2) {};
\node (v5) at (4, 1) {};
\node (v6) at (6, 1.5) {};
\node (v7) at (8, 2) {};
\node (v8) at (8, 1) {};

\begin{scope}[fill opacity = 0.7]

\filldraw [fill=lightorange] 
    ($(v1) + (-1,0)$)
    to [out=90, in=180] ($(v1) + (0,1)$)
    to [out=0, in=90] ($(v3) + (1,0)$)
    to [out=270, in=0] ($(v2) + (0,-1)$)
    to [out=180, in=270] ($(v1) + (-1,0)$);

\filldraw [fill=lightred] 
    ($(v4) + (1,0)$)
    to [out=90, in=0] ($(v4) + (0,1)$)
    to [out=180, in=90] ($(v3) + (-1,0)$)
    to [out=270, in=180] ($(v5) + (0,-1)$)
    to [out=0, in=270] ($(v4) + (1,0)$);

\filldraw [fill=lightpink] 
    ($(v4) + (-1,0)$)
    to [out=90, in=180] ($(v4) + (0,1)$)
    to [out=0, in=90] ($(v6) + (1,0)$)
    to [out=270, in=0] ($(v5) + (0,-1)$)
    to [out=180, in=270] ($(v4) + (-1,0)$);

\filldraw [fill=lightpurple] 
    ($(v7) + (1,0)$)
    to [out=90, in=0] ($(v7) + (0,1)$)
    to [out=180, in=90] ($(v6) + (-1,0)$)
    to [out=270, in=180] ($(v8) + (0,-1)$)
    to [out=0, in=270] ($(v7) + (1,0)$);

\end{scope}

\foreach \i in {1,...,8}
    \fill (v\i) circle (0.1);

\foreach \i/\name in {1/A[1],2/A[2],3/A[5],4/A[6],5/A[8],6/A[10],7/A[12],8/A[15]}
    \node[right] at (v\i) {$\name$};

\end{tikzpicture}

\caption{Hypergraph model of Figure~\ref{fig:csc592_bigdata_soe}. Each $A[i]$ is the $i$-th entry in the backing array $A$ (zero-indexed). Hyperedges 
correspond to keys $\mathcolorbox{lightorange}{x}$, $\mathcolorbox{lightred}{y}$, $\mathcolorbox{lightpink}{w}$ and $\mathcolorbox{lightpurple}{z}$, colored 
accordingly. Vertices of degree zero not shown.}
\label{fig:csc592_bigdata_hg}
\end{figure}

\noindent
In the filter setting, edges represent keys and vertices represent array entries. Higher edge density means lower space usage, since there are more keys (edges)
per array entry (vertex). Many families of random hypergraphs exhibit sharp phase transitions in peelability. When the edge density is below a critical value, 
called the 2-peelability threshold, sampled hypergraphs are almost surely 2-peelable. When the density exceeds this threshold, they are almost surely not. The 
threshold is therefore a direct measure of space efficiency. A higher threshold means more keys can be stored per array entry. \\

\noindent
Binary fuse filters induce the fuse graph family, which achieves the maximum 2-peelability threshold among 3-uniform hypergraph\footnote{In general, a 
hypergraph is $m$-uniform if every edge contains exactly $m$ vertices.} families. This threshold occurs at an edge density of $0.918$, or roughly 
$1.089 \lg \frac{1}{\epsilon}$ bits per key. No family of 3-uniform hypergraphs has higher 2-peelability threshold. More generally, a binary fuse filter with 
$m \geq 3$ hash functions\footnote{Only binary fuse filters with three or four hash functions are practical. With four hash functions, the 2-peelability 
threshold is already $0.977$, yielding about $1.023 \lg \frac{1}{\epsilon}$ bits per key or within $2.3\%$ of optimal.} saturates the 2-peelability threshold 
for $m$-uniform hypergraphs. \\

\noindent
However, this optimality result is confined to uniform hypergraphs. There may exist families of non-uniform hypergraphs with higher 2-peelability thresholds. If
such families also have a low average edge size (i.e. few hashes and accesses per query), filters based on them could potentially reduce memory overhead without
degrading query latency and throughput. Therefore, we ask:

\begin{enumerate}
    \item How can we systematically construct non-uniform hypergraph families with peelability thresholds comparable to or surpassing those of fuse graph 
    families?
    \item How can we design efficient approximate membership query data structures by exploiting peelable non-uniform hypergraph families?
\end{enumerate}

\subsection{Summary}

\subsubsection{Non-Uniform Hypergraphs}

We computed 2-peelability thresholds for 10K non-uniform fuse graph families with mean edge size between 3 and 4 (see 
Figure~\ref{fig:csc592_bigdata_nuhg_graph}). As peelability is monotonic in edge density, we used binary search. At each candidate threshold, we sampled 
three hypergraphs with 1M edges. If the majority were peelable, we considered the family peelable at that threshold. We implemented this algorithm in C++ and 
computed thresholds in parallel on a multicore cluster. \\

\begin{figure}[H]
    \centering
    \begin{subfigure}[b]{0.49\textwidth}
        \centering
        \begin{tikzpicture}
            \begin{axis}[
                xlabel={Mean Edge Size},
                ylabel={Peelability Threshold},
                xmin=3, xmax=4,
                ymin=0.88, ymax=0.94,
                width=\textwidth,
                legend to name=sharedlegend,
                legend columns=3,
                legend style={draw=none, column sep=4pt},
            ]
                \addplot[color=lightorange, mark=o, mark size=1pt, only marks, legend image post style={mark size=2pt, color=darkorange}] table [x=meanedge, y=threshold, col sep=comma] {data/csc592-bigdata-nu-1k.csv};
                \addlegendentry{Non-Uniform}
                \addplot[color=darkred, mark=*, mark size=1pt] table [x=meanedge, y=threshold, col sep=comma] {data/csc592-bigdata-nu-best.csv};
                \addlegendentry{Best Non-Uniform Sharded}
                \addplot[color=black!40, mark=*, mark size=1pt] table [x=meanedge, y=threshold, col sep=comma] {data/csc592-bigdata-shard.csv};
                \addlegendentry{Uniform Sharded}
            \end{axis}
        \end{tikzpicture}
    \end{subfigure}
    \hfill
    \begin{subfigure}[b]{0.49\textwidth}
        \centering
        \begin{tikzpicture}
            \begin{axis}[
                xlabel={Mean Hashes/Accesses},
                ylabel={Memory Overhead},
                yticklabel={\pgfmathprintnumber{\tick}\%},
                xmin=3, xmax=4,
                ymin=6, ymax=14,
                width=\textwidth
            ]
                \addplot[color=lightorange, mark=o, mark size=1pt, only marks] table [x=meanedge, y expr={100 * (1 / \thisrow{threshold} - 1)}, col sep=comma] {data/csc592-bigdata-nu-1k.csv};
                \addplot[color=darkred, mark=*, mark size=1pt] table [x=meanedge, y expr={100 * (1 / \thisrow{threshold} - 1)}, col sep=comma] {data/csc592-bigdata-nu-best.csv};
                \addplot[color=black!40, mark=*, mark size=1pt] table [x=meanedge, y expr={100 * (1 / \thisrow{threshold} - 1)}, col sep=comma] {data/csc592-bigdata-shard.csv};
            \end{axis}
        \end{tikzpicture}
    \end{subfigure}
    \ref{sharedlegend}
    \caption{Empirical peelability thresholds and memory overheads for non-uniform fuse graphs with 1M edges and mean edge size between 3 and 4. Orange dots 
    show 1K sampled families. Red shows the best non-uniform sharded families. Gray shows the sharded uniform baseline. Many non-uniform families outperform the
    uniform baseline. Left: mean edge size vs. peelability threshold (higher is better). Right: theoretical query cost vs. memory overhead (lower is better).}
    \label{fig:csc592_bigdata_nuhg_graph}
\end{figure}

\noindent
While $m$-uniform graphs have integer mean edge sizes, they can be extended to any mean edge size via sharding. For example, splitting edges evenly between 
edge-disjoint 3-uniform and 4-uniform subgraphs yields a mean edge size of 3.5. Many non-uniform families outperform this baseline, and therefore afford 
smaller filters. Interestingly, some non-uniform families also outperform uniform families at integer mean edge sizes. Since peelability thresholds are defined 
in the limit of large graphs, a non-uniform family can be more peelable than a uniform family with nearly the same mean edge size. For example, one family 
achieves $\approx 1$\% memory savings over the uniform baseline at mean edge size $4$. \\

\noindent
These data suggest two conjectures. Let $\mathcal{H}_d$ denote the set of fuse graph families with mean edge size $d$ and let $\tau(\mathcal{F})$ 
denote the 2-peelability threshold of a family $\mathcal{F}$. Define $\Gamma(d) = \sup \{ \tau(\mathcal{F}) : \mathcal{F} \in \mathcal{H}_d \}$. We conjecture 
that $\Gamma(d)$ is continuous, monotonically increasing, and concave. Since a concave $\Gamma$ would lie above any linear interpolation, this implies that 
sharding is suboptimal. We further conjecture that for $d \not \in \mathbb{Z}$, the supremum is attained by the family with a fraction $(1 - p)$ of 
$\lfloor d \rfloor$-edges and a fraction $p$ of $\lceil d \rceil$-edges, where $p = d - \lfloor d \rfloor$. For example, for $d = 3.5$, this family contains
50\% 3-edges and 50\% 4-edges. Both conjectures remain open.

\subsubsection{Non-Uniform Binary Fuse Filters}

\noindent
We implemented and optimized a non-uniform binary fuse filter with 50\% 3-edges and 50\% 4-edges in C++ (i.e. each key hashes to either 3 or 4 array entries). 
This edge distribution offered the greatest advantage over sharding uniform binary fuse filters. We considered three query algorithms: conditional branch, 
branchless and phantom entry (see Table~\ref{tab:csc592_bigdata_impl_tab} for metrics and see Section~\ref{sec:csc592_bigdata_contains_impl} for pseudocode). \\

\begin{table}[H]
    \centering
    \begin{tabular}{lccc}
        \toprule
        Algorithm & Hashes & Branch Mispredictions & Cache Misses \\
        \midrule
        Conditional Branch & 3.5 & 0.5 & 3.5 \\
        Branchless        & 4   & 0   & 4   \\
        Phantom Entry     & 4   & 0   & 3.5 \\
    \bottomrule
    \end{tabular}
    \caption{Metrics for non-uniform binary fuse filter query algorithms. Hashes, branch mispredictions, and cache misses are mean per query.}
    \label{tab:csc592_bigdata_impl_tab}
\end{table}

\noindent
In the conditional branch algorithm, we compute the fingerprint and access the first three array entries, then compute their bitwise-XOR. If the key is a 
4-edge, we also access and bitwise-XOR the fourth array entry. This branch is mispredicted for at least half of the queries. \\

\noindent
In the branchless algorithm, we compute the fingerprint and access all four array entries. If the key is a 3-edge, we zero the fourth value with a 
bitwise-AND. Then, we compute the bitwise-XOR of all five values. For large filters, all four array accesses will miss the cache. \\ 

\noindent
In the phantom entry algorithm, we compute the fingerprint and access all four array entries. If the key is a 3-edge, we zero the fourth hash with a bitwise-AND
before accessing its entry. That is, the fourth array entry for all 3-edges is $A[0]$. We shift the original entries left once and set $A[0]$ to zero (see 
Figure~\ref{fig:csc592_bigdata_nubff}). Unlike the branchless algorithm, the fourth array access for 3-edges will hit the cache since $A[0]$ is requested in 
50\% of queries. \\

\begin{figure}[H]
    \centering
    \begin{tikzpicture}[>=Stealth, arraynode/.style={draw, minimum width=7mm, minimum height=6mm, outer sep=0, inner sep=0}]

        \node[arraynode, fill=lightgray] (a1) {000};
        \node[arraynode, fill=lightorange, right=0 of a1] (a2) {001};
        \node[arraynode, fill=lightyellow, right=0 of a2] (a3) {000};
        \node[arraynode, right=0 of a3, path picture={
        \shade[left color=lightorange, right color=lightred]
                (path picture bounding box.south west)
                rectangle (path picture bounding box.north east);
        }] (a4) {101};
        \node[arraynode, fill=lightyellow, right=0 of a4] (a5) {000};
        \node[arraynode, right=0 of a5, path picture={
        \shade[left color=lightorange, right color=lightred]
                (path picture bounding box.south west)
                rectangle (path picture bounding box.north east);
        }] (a6) {000};
        \node[arraynode, fill=lightyellow, right=0 of a6] (a7) {000};
        \node[arraynode, fill=lightorange, right=0 of a7] (a8) {000};
        \node[arraynode, fill=lightyellow, right=0 of a8, path picture={
        \shade[left color=lightred, right color=lightpink]
                (path picture bounding box.south west)
                rectangle (path picture bounding box.north east);
        }] (a9) {110};
        \node[arraynode, fill=lightpink, right=0 of a9] (a10) {000};
        \node[arraynode, fill=lightpurple, right=0 of a10] (a11) {111};
        \node[arraynode, fill=lightyellow, right=0 of a11] (a12) {000};
        \node[arraynode, right=0 of a12, path picture={
        \shade[left color=lightred, right color=lightpink]
                (path picture bounding box.south west)
                rectangle (path picture bounding box.north east);
        }] (a13) {000};
        \node[arraynode, fill=lightpink, right=0 of a13] (a14) {000};
        \node[arraynode, fill=lightpurple, right=0 of a14] (a15) {000};

        \draw[thick] 
            ([yshift=4pt]a1.north west) -- 
            ([yshift=-4pt]a1.south west);
        
        \foreach \i in {1,3,5,7,9,11,13,15} {
            \draw[thick] 
                ([yshift=4pt]a\i.north east) -- 
                ([yshift=-4pt]a\i.south east);
        }

        \node[circle, draw, above=12mm of a4, fill=lightorange] (x) {$x$};
        \draw[->] (x) to (a2.north);
        \draw[->] (x) to (a4.north);
        \draw[->] (x) to (a6.north);
        \draw[->] (x) to (a8.north);
        
        \node[circle, draw, above=12mm of a7, fill=lightred] (y) {$y$};
        \draw[->] (y) to (a4.north);
        \draw[->] (y) to (a6.north);
        \draw[->] (y) to (a9.north);
        \draw[->, dashed, lightgray] (y) to[out=150, in=110] (a1.north);

        \node[circle, draw, above=12mm of a10, fill=lightpink] (w) {$w$};
        \draw[->] (w) to (a9.north);
        \draw[->] (w) to (a10.north);
        \draw[->] (w) to (a13.north);
        \draw[->] (w) to (a14.north);

        \node[circle, draw, above=12mm of a13, fill=lightpurple] (z) {$z$};
        \draw[->] (z) to (a11.north);
        \draw[->] (z) to (a13.north);
        \draw[->] (z) to (a15.north);
        \draw[->, dashed, lightgray] (z) to[out=160, in=110] (a1.north);
        
    \end{tikzpicture}
    \caption{Non-uniform binary fuse filter with phantom entry, 14-entry array, 2-entry segments, and 3-bit fingerprints. Vertical bars mark segment boundaries.
    The phantom entry is colored gray. The filter contains keys $\mathcolorbox{lightorange}{x}$, $\mathcolorbox{lightred}{y}$, $\mathcolorbox{lightpink}{w}$ and 
    $\mathcolorbox{lightpurple}{z}$, with fingerprints $\mathcolorbox{lightorange}{100}$, $\mathcolorbox{lightred}{011}$, $\mathcolorbox{lightpink}{110}$ and 
    $\mathcolorbox{lightpurple}{111}$, respectively. Entries hashed to by keys are marked with arrows and the keys' colors. Each key hashes to three or four 
    consecutive segments. Dashed gray arrows indicate hashes to the phantom entry.}
    \label{fig:csc592_bigdata_nubff}
\end{figure}


\noindent
Surprisingly, only the conditional branch algorithm had better query throughput than the 4-uniform binary fuse filter. We expected that nearly all cycles were
spent stalled on loads, so a lower cache miss rate should improve throughput by generating useful outstanding misses. However, only $80\%$ of cycles 
were spent stalled (all loads), not nearly 100\%. We hypothesize that either the ROB is too small to allow OoO address generation before retiring outstanding 
misses, or lack of MSHRs stalls OoO address generation when loading the next key. This means hash latency is the primary bottleneck.

\subsection{Reflection}

TODO: A 250-word self reflection of what skills were developed or enhanced in completing those
projects.